% GENERATED FILE. Edit canonical lessons, then run npm run generate:books.
% canonical-source-hash: fnv1a64:8d1eab0e70734066
% canonical-lessons: BN-C262-lekhok, BN-C262-kobi, BN-C262-shilpi, BN-C262-gayika, BN-C262-nortoki, BN-R262-first-pass-words-from-chapters-255-258, BN-R262-second-pass-words-from-chapters-259-262

\chapter{Writer, Poet, Artist, Female singer, Female dancer}
\label{ch:bn-writer-poet-artist-female-singer-female-dancer}
\hlchaptermodality{262}

\begin{chapteropening}
\textbf{By the end of this chapter:} \emph{I can say a writer, a poet, an artist, a female singer and a female dancer.}

Complete the last lesson of chapter 262: I can say a writer, a poet, an artist, a female singer and a female dancer.
\end{chapteropening}

\section[\texorpdfstring{\bn{লেখক} (lekhôk)}{লেখক (lekhôk)}]{\bn{লেখক} (lekhôk) — a writer}

\label{lesson:BN-C262-lekhok}



\begin{quote}
Before the new one: say the Bengali for a gatekeeper, then the Bengali for an army.
\end{quote}

\subsection*{You'll want to know: \bn{লেখক}}
\textbf{\bn{লেখক}} — \emph{lekhôk} — \textquotedblleft{}a writer\textquotedblright{}.

\begin{grammarlens}[title={What you've built}]
One more word for this chapter.
\end{grammarlens}

\subsection*{Your turn}
\begin{itemize}
\raggedright
  \item \emph{Say it:} \emph{lekhôk}
  \item \emph{Say it:} \emph{lekhôk}, once more
  \item \emph{Say it:} say \emph{lekhôk}
  \item \emph{Recall:} say the Bengali for a gatekeeper, then the Bengali for an army, then say \emph{lekhôk} again
\end{itemize}

\begin{tcolorbox}[breakable,colback=teal!4,colframe=teal!35!black,title={Before you move on}]
What does \bn{লেখক} mean? (\textquotedblleft{}A writer\textquotedblright{}.) Say it once more.
\end{tcolorbox}
\section[\texorpdfstring{\bn{কবি} (kôbi)}{কবি (kôbi)}]{\bn{কবি} (kôbi) — a poet}

\label{lesson:BN-C262-kobi}



\begin{quote}
Before the new one: say the Bengali for an army, then the Bengali for a writer.
\end{quote}

\subsection*{You'll want to know: \bn{কবি}}
\textbf{\bn{কবি}} — \emph{kôbi} — \textquotedblleft{}a poet\textquotedblright{}.

\begin{grammarlens}[title={What you've built}]
One more word for this chapter.
\end{grammarlens}

\subsection*{Your turn}
\begin{itemize}
\raggedright
  \item \emph{Say it:} \emph{kôbi}
  \item \emph{Say it:} \emph{kôbi}, once more
  \item \emph{Say it:} say \emph{kôbi}
  \item \emph{Recall:} say the Bengali for an army, then the Bengali for a writer, then say \emph{kôbi} again
\end{itemize}

\begin{tcolorbox}[breakable,colback=teal!4,colframe=teal!35!black,title={Before you move on}]
What does \bn{কবি} mean? (\textquotedblleft{}A poet\textquotedblright{}.) Say it once more.
\end{tcolorbox}
\section[\texorpdfstring{\bn{শিল্পী} (shilpī)}{শিল্পী (shilpī)}]{\bn{শিল্পী} (shilpī) — an artist}

\label{lesson:BN-C262-shilpi}



\begin{quote}
Before the new one: say the Bengali for a writer, then the Bengali for a poet.
\end{quote}

\subsection*{You'll want to know: \bn{শিল্পী}}
\textbf{\bn{শিল্পী}} — \emph{shilpī} — \textquotedblleft{}an artist\textquotedblright{}.

\begin{grammarlens}[title={What you've built}]
One more word for this chapter.
\end{grammarlens}

\subsection*{Your turn}
\begin{itemize}
\raggedright
  \item \emph{Say it:} \emph{shilpī}
  \item \emph{Say it:} \emph{shilpī}, once more
  \item \emph{Say it:} say \emph{shilpī}
  \item \emph{Recall:} say the Bengali for a writer, then the Bengali for a poet, then say \emph{shilpī} again
\end{itemize}

\begin{tcolorbox}[breakable,colback=teal!4,colframe=teal!35!black,title={Before you move on}]
What does \bn{শিল্পী} mean? (\textquotedblleft{}An artist\textquotedblright{}.) Say it once more.
\end{tcolorbox}
\section[\texorpdfstring{\bn{গায়িকা} (gāyikā)}{গায়িকা (gāyikā)}]{\bn{গায়িকা} (gāyikā) — a female singer}

\label{lesson:BN-C262-gayika}



\begin{quote}
Before the new one: say the Bengali for a poet, then the Bengali for an artist.
\end{quote}

\subsection*{You'll want to know: \bn{গায়িকা}}
\textbf{\bn{গায়িকা}} — \emph{gāyikā} — \textquotedblleft{}a female singer\textquotedblright{}.

\begin{grammarlens}[title={What you've built}]
One more word for this chapter.
\end{grammarlens}

\subsection*{Your turn}
\begin{itemize}
\raggedright
  \item \emph{Say it:} \emph{gāyikā}
  \item \emph{Say it:} \emph{gāyikā}, once more
  \item \emph{Say it:} say \emph{gāyikā}
  \item \emph{Recall:} say the Bengali for a poet, then the Bengali for an artist, then say \emph{gāyikā} again
\end{itemize}

\begin{tcolorbox}[breakable,colback=teal!4,colframe=teal!35!black,title={Before you move on}]
What does \bn{গায়িকা} mean? (\textquotedblleft{}A female singer\textquotedblright{}.) Say it once more.
\end{tcolorbox}
\section[\texorpdfstring{\bn{নর্তকী} (nôrtôkī)}{নর্তকী (nôrtôkī)}]{\bn{নর্তকী} (nôrtôkī) — a female dancer}

\label{lesson:BN-C262-nortoki}



\begin{quote}
Before the new one: say the Bengali for an artist, then the Bengali for a female singer.
\end{quote}

\subsection*{You'll want to know: \bn{নর্তকী}}
\textbf{\bn{নর্তকী}} — \emph{nôrtôkī} — \textquotedblleft{}a female dancer\textquotedblright{}.

\begin{grammarlens}[title={What you've built}]
One more word for this chapter.
\end{grammarlens}

\subsection*{Your turn}
\begin{itemize}
\raggedright
  \item \emph{Say it:} \emph{nôrtôkī}
  \item \emph{Say it:} \emph{nôrtôkī}, once more
  \item \emph{Say it:} say \emph{nôrtôkī}
  \item \emph{Recall:} say the Bengali for an artist, then the Bengali for a female singer, then say \emph{nôrtôkī} again
\end{itemize}

\begin{tcolorbox}[breakable,colback=teal!4,colframe=teal!35!black,title={Before you move on}]
What does \bn{নর্তকী} mean? (\textquotedblleft{}A female dancer\textquotedblright{}.) Say it once more.
\end{tcolorbox}
\section[(dialogue)]{Words from chapters 255-258 — a first pass}

\label{lesson:BN-R262-first-pass-words-from-chapters-255-258}



\begin{quote}
Say the words for a female singer and a female dancer.
\end{quote}

\subsection*{Your turn}
Say these aloud:

\begin{itemize}
\raggedright
  \item sunshine, a storm, a flood, smoke, ash
  \item culture, history, science, art; industry, literature
  \item chhana, paneer, ice cream, chira, phuchka
  \item an email, a message, a password, a battery, a switch
\end{itemize}

\begin{tcolorbox}[breakable,colback=teal!4,colframe=teal!35!black,title={Before you move on}]
Sunshine? (\textbf{\bn{রোদ}}, \emph{rod}.) A storm? (\textbf{\bn{ঝড়}}, \emph{jhôṛ}.) A flood? (\textbf{\bn{বন্যা}}, \emph{bônyā}.) Smoke? (\textbf{\bn{ধোঁয়া}}, \emph{dhõyā}.) Ash? (\textbf{\bn{ছাই}}, \emph{chhāi}.) Culture? (\textbf{\bn{সংস্কৃতি}}, \emph{sôngskriti}.) History? (\textbf{\bn{ইতিহাস}}, \emph{itihās}.) Science? (\textbf{\bn{বিজ্ঞান}}, \emph{bigyān}.) Art; industry? (\textbf{\bn{শিল্প}}, \emph{shilpô}.) Literature? (\textbf{\bn{সাহিত্য}}, \emph{sāhityô}.) Chhana? (\textbf{\bn{ছানা}}, \emph{chhānā}.) Paneer? (\textbf{\bn{পনির}}, \emph{pônir}.) Ice cream? (\textbf{\bn{আইসক্রিম}}, \emph{āiskrim}.) Chira? (\textbf{\bn{চিঁড়ে}}, \emph{chĩṛe}.) Phuchka? (\textbf{\bn{ফুচকা}}, \emph{phuchkā}.) An email? (\textbf{\bn{ইমেল}}, \emph{imel}.) A message? (\textbf{\bn{বার্তা}}, \emph{bārtā}.) A password? (\textbf{\bn{পাসওয়ার্ড}}, \emph{pāsoyārḍ}.) A battery? (\textbf{\bn{ব্যাটারি}}, \emph{bæṭāri}.) A switch? (\textbf{\bn{সুইচ}}, \emph{suich}.)
\end{tcolorbox}
\section[(dialogue)]{Words from chapters 259-262 — a second pass}

\label{lesson:BN-R262-second-pass-words-from-chapters-259-262}



\begin{quote}
Say the words for a rate and a discount.
\end{quote}

\subsection*{Your turn}
Say these aloud:

\begin{itemize}
\raggedright
  \item a rate, a discount, a form, an application, an identity card
  \item a longing, a prize, a punishment, a competition, victory
  \item a mechanic, a cook, a waiter, a gatekeeper, an army
  \item a writer, a poet, an artist, a female singer, a female dancer
\end{itemize}

\begin{tcolorbox}[breakable,colback=teal!4,colframe=teal!35!black,title={Before you move on}]
A rate? (\textbf{\bn{দর}}, \emph{dôr}.) A discount? (\textbf{\bn{ছাড়}}, \emph{chhāṛ}.) A form? (\textbf{\bn{ফর্ম}}, \emph{phôrm}.) An application? (\textbf{\bn{আবেদন}}, \emph{ābedôn}.) An identity card? (\textbf{\bn{পরিচয়পত্র}}, \emph{pôrichôypôtrô}.) A longing? (\textbf{\bn{সাধ}}, \emph{sādh}.) A prize? (\textbf{\bn{পুরস্কার}}, \emph{purôskār}.) A punishment? (\textbf{\bn{শাস্তি}}, \emph{shāsti}.) A competition? (\textbf{\bn{প্রতিযোগিতা}}, \emph{prôtijogitā}.) Victory? (\textbf{\bn{জয়}}, \emph{jôy}.) A mechanic? (\textbf{\bn{মিস্ত্রি}}, \emph{mistri}.) A cook? (\textbf{\bn{রাঁধুনি}}, \emph{rā̃dhuni}.) A waiter? (\textbf{\bn{ওয়েটার}}, \emph{oyeṭār}.) A gatekeeper? (\textbf{\bn{দারোয়ান}}, \emph{dāroyān}.) An army? (\textbf{\bn{সেনাবাহিনী}}, \emph{senābāhinī}.) A writer? (\textbf{\bn{লেখক}}, \emph{lekhôk}.) A poet? (\textbf{\bn{কবি}}, \emph{kôbi}.) An artist? (\textbf{\bn{শিল্পী}}, \emph{shilpī}.) A female singer? (\textbf{\bn{গায়িকা}}, \emph{gāyikā}.) A female dancer? (\textbf{\bn{নর্তকী}}, \emph{nôrtôkī}.)
\end{tcolorbox}
